%% 03_metodologia.tex — Testbed, oráculo, diseño RCBD, reproducibilidad.

\section{Metodología experimental}
\label{sec:metodologia}

\subsection{Banco de pruebas}
\label{sec:testbed}

El objetivo de evaluación es OWASP Benchmark~\cite{owaspbenchmark}, una
aplicación Java deliberadamente vulnerable con casos de prueba etiquetados. Se
despliega mediante \code{docker compose}, sirviendo HTTPS en el puerto
\num{8443} con un certificado autofirmado (de ahí que todas las sondas usen
\code{-{}-no-verify-ssl}). Para acotar el estudio a un único tipo de vector, el
banco se restringe a los \textbf{\num{338}~casos de parámetro \code{GET}}
(\code{vector = getparam}). El generador \code{build\_ground\_truth.py} produce
el fichero de verdad-terreno como un arreglo de registros
$\langle \code{url}, \code{param}, \code{type}, \code{vulnerable} \rangle$; los
registros con \code{vulnerable = false} son \emph{trampas} explícitas (puntos
conocidos como seguros que no deben dispararse). La \cref{tab:testbed} resume
su composición.

\begin{table}[t]
  \centering
  \caption{Composición de la verdad-terreno (OWASP Benchmark, vectores
  \code{getparam}).}
  \label{tab:testbed}
  \begin{tabular}{lrr}
    \toprule
    Categoría & Instancias vulnerables & Trampas \\
    \midrule
    XSS (reflejado)          & 68 & \multirow{5}{*}{144} \\
    Inyección SQL            & 65 & \\
    Path Traversal           & 30 & \\
    Inyección de comandos    & 25 & \\
    Inyección LDAP           & 6  & \\
    \midrule
    \textbf{Total}           & \textbf{194} & \textbf{144} \\
    \bottomrule
  \end{tabular}
\end{table}

\subsection{Oráculo de calificación}
\label{sec:oraculo}

La calificación la realiza \code{tools/eval\_oracle.py}, un guion independiente
sin dependencias de terceros. Normaliza cada tupla
$(\code{url}, \code{param}, \code{type})$ con las mismas funciones de
\emph{folding} que emplea el análisis (se descartan cadena de consulta y
fragmento; los alias de tipo se colapsan sobre un mismo \emph{bucket}), y
calcula:
\begin{align}
  \mathrm{TP} &= \lvert H \cap P \rvert, &
  \mathrm{FP} &= \lvert H \setminus P \rvert, \nonumber\\
  \mathrm{FN} &= \lvert P \setminus H \rvert, &
  \mathrm{FPR} &= \frac{\lvert H \cap T \rvert}{\lvert T \rvert},
  \label{eq:metricas}
\end{align}
donde $H$ es el conjunto de hallazgos, $P$ el de instancias vulnerables y $T$
el de trampas. La FPR se define \emph{exclusivamente} sobre la población de
trampas conocidas; cuando la verdad-terreno no define trampas, se reporta
\code{N/A}. De aquí se derivan precisión, exhaustividad (\emph{recall}) y
$F_1$.

\subsection{Diseño experimental}
\label{sec:diseno}

Se emplea un \textbf{diseño de bloques completos al azar} (RCBD) en el que:

\begin{itemize}
  \item El \textbf{factor} es la configuración del \emph{pipeline}, con cuatro
  niveles: \code{baseline} (todas las optimizaciones activas),
  \code{no-interleave}, \code{no-priority} y \code{no-runtime-confirm} --- es
  decir, una ablación de un solo componente por brazo.
  \item El \textbf{presupuesto de peticiones}
  (\code{-{}-max-payloads} $\in \{10, 20, 50, \infty\}$) se trata como un
  factor de bloqueo: la detección es trivialmente monótona en el presupuesto,
  y bloquear por él aísla el efecto de la configuración.
  \item La \textbf{unidad de bloqueo} para las pruebas de significancia es el
  \emph{coste de detección por instancia vulnerable}: el \code{request\_index}
  de la primera sonda que la detecta, censurado en $(\text{presupuesto}+1)$
  cuando nunca se detecta.
\end{itemize}

La rejilla completa es $4 \times 4 = \NumCorridas{}$ ejecuciones, orquestadas
por \code{tools/run\_experiments.py} y consolidadas en
\code{testbed/experiment\_results.csv}; las \NumCorridas{} se calificaron sin
excepción, de modo que los \NumBloquesCompletos{} bloques de presupuesto
quedan completos para las pruebas de significancia (\cref{sec:resultados}).
Cada ejecución produce su propia telemetría JSONL, su informe JSON y su
\code{metrics.json}.

\subsection{Métrica de eficiencia: AUC de la curva de detección}
\label{sec:auc}

Para cada ejecución se construye la curva escalonada
$c(k) = \lvert \{\text{instancias detectadas con } \code{request\_index} \le k\} \rvert$.
Tomada sobre los ejes normalizados (exhaustividad frente a fracción de
presupuesto), su área
\begin{equation}
  \mathrm{AUC} = \int_0^1 \frac{c(k)}{\lvert P \rvert}\, \mathrm{d}\hat{k}
  \label{eq:auc}
\end{equation}
equivale a la exhaustividad media sobre el presupuesto y es la medida principal
de eficiencia: un \emph{pipeline} que detecta lo mismo pero antes tiene mayor
AUC.

\subsection{Pruebas estadísticas}
\label{sec:tests}

Para cada presupuesto en el que las cuatro configuraciones completaron su
ejecución se comparan los cuatro vectores de coste de detección con:
\begin{enumerate}
  \item una prueba de \textbf{Friedman}~\cite{friedman1937use} sobre las cuatro
  condiciones (bloques $=$ instancias vulnerables);
  \item un post-hoc de \textbf{Nemenyi}~\cite{nemenyi1963distribution} para las
  comparaciones por pares;
  \item \textbf{Wilcoxon} de rangos con signo y \textbf{McNemar} para
  \code{baseline} frente a cada ablación, con corrección de
  \textbf{Benjamini--Hochberg}~\cite{benjamini1995controlling} sobre la familia
  de pares;
  \item $\delta$ de \textbf{Cliff}~\cite{cliff1993dominance} como tamaño del
  efecto no paramétrico (un $\delta > 0$ indica que la ablación es más lenta o
  peor).
\end{enumerate}
Este protocolo sigue las recomendaciones de Dem\v{s}ar para la comparación de
métodos sobre múltiples conjuntos~\cite{demsar2006statistical}. Los bloques se
agrupan además sobre todos los presupuestos (\emph{pooled}) para una prueba
global.

\subsection{Pilares de reproducibilidad}
\label{sec:reprod}

Toda ejecución reportada va acompañada de: (i)~el manifiesto de la
\cref{sec:manifiesto}, con el \emph{hash} SHA-256 del corpus y el \emph{commit}
de \code{git}; (ii)~una semilla global (\code{-{}-global-seed}) que fija la
ordenación \code{random} de \emph{payloads}, la rotación de \emph{User-Agent},
el \emph{jitter} y la mutación de caja; y (iii)~el guion de orquestación, que
es idempotente --- un directorio de resultados existente se omite salvo
\code{-{}-force}, de modo que una barrida interrumpida se reanuda sin
recomputar.
